\section{命名战争与开源扩散：产品之外的执行力}

OpenClaw 的传播史里，一个看似琐碎却关键的事件是多轮 rename。从 Claudebot/Codex 到 OpenClaw，背后是品牌冲突、生态关系与社区协同。

\subsection{Rename 为什么会成为工程任务}

访谈里提到，“仅 Codex 到 OpenClaw 的改名就花了约 10 小时”，因为不是字符串替换，而是仓库内外、域名、包名、账号引用的系统迁移。

\begin{knowledgebox}{Rename 的隐性复杂度}
\begin{itemize}
    \item 包管理生态（依赖名、缓存、镜像）；
    \item 文档生态（教程、截图、外部引用）；
    \item 社交生态（账号名、组织名、历史链接）；
    \item 法务生态（商标近似、品牌混淆风险）。
\end{itemize}
\end{knowledgebox}

\begin{figure}[H]
\centering
\includegraphics[width=0.9\textwidth]{frames/frame-003523.jpg}
\caption{改名战阶段：品牌冲突、账号迁移与社区协调并行\protect\footnotemark}
\end{figure}
\footnotetext{视频画面时间区间：00:35:15--00:35:30。}

\subsection{“叫得出来”是开源项目的重要基础设施}

一个项目若无法被稳定命名、检索、引用，其扩散会被严重折损。OpenClaw 的案例说明，开源维护者需要同时具备技术执行力与传播执行力。

\begin{importantbox}{从改名事件看项目管理能力}
Peter 在该阶段展示的不是“营销技巧”，而是高压下的组织能力：快速决策、并行协作、对外沟通、对内一致性推进。
\end{importantbox}

\subsection{本章小结}
OpenClaw 的扩散不只靠“功能新奇”，还靠运营与工程并重的执行体系。对开源项目而言，命名、文档、迁移同样属于核心工程。

